\documentclass{beamer}
\usepackage[utf8]{inputenc}

\usetheme{Madrid}
\usecolortheme{default}
\usepackage{amsmath,amssymb,amsfonts,amsthm}
\usepackage{txfonts}
\usepackage{tkz-euclide}
\usepackage{listings}
\usepackage{adjustbox}
\usepackage{array}
\usepackage{tabularx}
\usepackage{gvv}
\usepackage{lmodern}
\usepackage{circuitikz}
\usepackage{tikz}
\usepackage[utf8]{inputenc}

\usetheme{Madrid}
\usecolortheme{default}
\usepackage{amsmath,amssymb,amsfonts,amsthm}
\usepackage{txfonts}
\usepackage{tkz-euclide}
\usepackage{listings}
\usepackage{adjustbox}
\usepackage{array}
\usepackage{tabularx}
\usepackage{gvv}
\usepackage{lmodern}
\usepackage{circuitikz}
\usepackage{tikz}
\usepackage{graphicx}

\setbeamertemplate{page number in head/foot}[totalframenumber]

\usepackage{tcolorbox}
\tcbuselibrary{minted,breakable,xparse,skins}



\definecolor{bg}{gray}{0.95}
\DeclareTCBListing{mintedbox}{O{}m!O{}}{%
  breakable=true,
  listing engine=minted,
  listing only,
  minted language=#2,
  minted style=default,
  minted options={%
    linenos,
    gobble=0,
    breaklines=true,
    breakafter=,,
    fontsize=\small,
    numbersep=8pt,
    #1},
  boxsep=0pt,
  left skip=0pt,
  right skip=0pt,
  left=25pt,
  right=0pt,
  top=3pt,
  bottom=3pt,
  arc=5pt,
  leftrule=0pt,
  rightrule=0pt,
  bottomrule=2pt,

  colback=bg,
  colframe=orange!70,
  enhanced,
  overlay={%
    \begin{tcbclipinterior}
    \fill[orange!20!white] (frame.south west) rectangle ([xshift=20pt]frame.north west);
    \end{tcbclipinterior}},
  #3,
}
\lstset{
    language=C,
    basicstyle=\ttfamily\small,
    keywordstyle=\color{blue},
    stringstyle=\color{orange},
    commentstyle=\color{green!60!black},
    numbers=left,
    numberstyle=\tiny\color{gray},
    breaklines=true,
    showstringspaces=false,
}
%------------------------------------------------------------
%This block of code defines the information to appear in the
%Title page
\title %optional
{4.10.16}
\date{September  2025}
%\subtitle{A short story}

\author % (optional)
{Bhoomika L - EE25BTECH11014}

\begin{document}


\frame{\titlepage}
\begin{frame}{Question}
The area of the region bounded by the curve y = x + 1 and the lines x = 2
and x = 3 is:
\begin{multicols}
\begin{enumerate}
    \item $\frac{7}{2} sq.units$
    \item $\frac{9}{2} sq.units$
    \item $\frac{11}{2} sq.units$
    \item $\frac{13}{2} sq.units$
\end{enumerate}
\end{multicols}
\end{frame}

\begin{frame}{Solution}
y = x + 1 can be expressed as
\begin{align}
 \vec{n^Tx}=c\\
\vec{n^T}=\myvec{1\\-1} \& c=-1
\end{align}

where $\vec{n^T}$ is the normal vector of the line.\\
Let the points of intersection of the curve $y = x + 1$ and the lines $x=2$ and $x=3$ be
$\vec{A}$ and $\vec{B}.$
\end{frame}

\begin{frame}{Solution}
For $\vec{A}$:
\begin{align}
\myvec{1 & 0 \\ 
-1 & 1}
\myvec{x\\y}=\myvec{2\\1}
\end{align}
The augmented matrix of the above is:
\begin{align}
\myvec{\begin{array}{cc|c}
1 & 0 & 2 \\ 
-1 & 1 & 1
\end{array}}.
\end{align}
Then, 
\begin{align}
    \vec{R_1}\longrightarrow\vec{R_2}+\vec{R_1}\\
  \myvec{\begin{array}{cc|c}
1 & 0 & 2 \\ 
0 & 1 & 3
\end{array}}\\
\therefore \myvec{x\\y}=\myvec{2\\3}.
\end{align}    
\end{frame}

\begin{frame}{Solution}
For $\vec{B}$:
\begin{align}
\myvec{1 & 0 \\ 
-1 & 1}
\myvec{x\\y}=\myvec{3\\1}
\end{align}

The augmented matrix of the above is:
\begin{align}
\myvec{\begin{array}{cc|c}
1 & 0 & 3 \\ 
-1 & 1 & 1
\end{array}}.
\end{align}

Then, 
\begin{align}
    R_2 \longrightarrow R_2 + R_1 \\
  \myvec{\begin{array}{cc|c}
1 & 0 & 3 \\ 
0 & 1 & 4
\end{array}}\\
\therefore  \myvec{x\\y}=\myvec{3\\4}.
\end{align}
\end{frame}

\begin{frame}{Solution}
Hence the points of the quadrilateral are $\brak{2,0},\brak{2,3},\brak{3,4}$ and $\brak{3,0}$ (which are $\vec{A},\vec{B},\vec{C}$ and $\vec{D}$ respectively.)

 Area of the quadrilateral can be calculated by using the formula,
 \begin{align}
 Ar=\frac{1}{2}\left\|\brak{\vec{A}-\vec{C}}\times\brak{\vec{D}-\vec{B}}\right\|\\
 \brak{\vec{A}-\vec{C}}=\myvec{-1\\-4}   \\\brak{\vec{B}-\vec{D}}=\myvec{1\\3}  \\
 \brak{\vec{A}-\vec{C}}\times\brak{\vec{D}-\vec{B}}=\begin{vmatrix}-1&1 \\-4&3\end{vmatrix}=-3-4=-7\\
 \therefore Ar=\frac{1}{2}\left|-7\right|=\frac{7}{2}.
 \end{align}
 Therefore (option 1) $\frac{7}{2} sq.units$ is the correct answer.
\end{frame}

\begin{frame}{Plot}
    \centering
    \includegraphics[width=\columnwidth, height=0.8\textheight, keepaspectratio]{../figs/fig_8.png}     
\end{frame}

 \begin{frame}[fragile]
\frametitle{C Code }
\begin{lstlisting}

// area.c
// Function to compute area under y = x + 1 from a to b

double area_under_line(double a, double b) {
    if (a > b) {  // swap if needed
        double t = a;
        a = b;
        b = t;
    }
    return 0.5 * (b*b - a*a) + (b - a);
}
\end{lstlisting}
\end{frame}

 \begin{frame}[fragile]
\frametitle{Python Code}
\begin{lstlisting}
import numpy as np
import matplotlib.pyplot as plt
from ctypes import CDLL, c_double

lib = CDLL("./area.so")

# Set argument and return types
lib.area_under_line.argtypes = (c_double, c_double)
lib.area_under_line.restype = c_double

# Example usage
a, b = 2.0, 3.0
area = lib.area_under_line(a, b)

print(f"Area under y = x + 1 from {a} to {b} is {area:.6f}")

\end{lstlisting}
\end{frame}

 \begin{frame}[fragile]
\frametitle{Python Code}
\begin{lstlisting}
# Define the function
def f(x):
    return x + 1

# Range for the curve
x_vals = np.linspace(1.5, 3.5, 200)
y_vals = f(x_vals)

# Intersection points
points = {
    "A": (2, f(2)),   # (2,3)
    "B": (3, f(3)),   # (3,4)
    "C": (3, 0),      # (3,0)
    "D": (2, 0)       # (2,0)
}
\end{lstlisting}
\end{frame}

 \begin{frame}[fragile]
\frametitle{Python Code}
\begin{lstlisting}
# Plot the line y = x+1
plt.plot(x_vals, y_vals, label=r"$y=x+1$", color="blue")

# Plot vertical lines x=2 and x=3
plt.axvline(2, color="black", linestyle="--", label=r"$x=2$")
plt.axvline(3, color="black", linestyle="--", label=r"$x=3$")

# Plot x-axis
plt.axhline(0, color="black")

\end{lstlisting}
\end{frame}

 \begin{frame}[fragile]
\frametitle{Python Code}
\begin{lstlisting}
# Shade the region (quadrilateral A-B-C-D)
x_poly = [points["A"][0], points["B"][0], points["C"][0], points["D"][0]]
y_poly = [points["A"][1], points["B"][1], points["C"][1], points["D"][1]]
plt.fill(x_poly, y_poly, color="skyblue", alpha=0.5, label="Shaded Region")

# Mark the intersection points
for name, (x, y) in points.items():
    plt.scatter(x, y, color="red")
    plt.text(x+0.05, y+0.1, f"{name}({x},{y})")

\end{lstlisting}
\end{frame}

 \begin{frame}[fragile]
\frametitle{Python Code}
\begin{lstlisting}
# Labels and legend
plt.xlabel("x-axis")
plt.ylabel("y-axis")
plt.title("Region bounded by y=x+1, x=2, x=3 and x-axis")
plt.legend()
plt.grid(True)
plt.axis("equal")

# Save the figure
plt.savefig("fig_8.png", dpi=300, bbox_inches="tight")

# Show the plot
plt.show()
\end{lstlisting}
\end{frame}

\end{document}
